% ============================================================
%  AGROSIGNAL — Описание набора данных и моделей
%  ПВКР.ОД.001-2026
%  Компиляция: xelatex README_DATA.tex  (дважды для оглавления)
% ============================================================
\documentclass[14pt,a4paper]{extarticle}

%% ── Шрифты ──────────────────────────────────────────────────────────────────
\usepackage{fontspec}
\setmainfont[
  BoldFont       = {Times New Roman Bold},
  ItalicFont     = {Times New Roman Italic},
  BoldItalicFont = {Times New Roman Bold Italic},
  Ligatures      = TeX
]{Times New Roman}
\setmonofont[Scale=0.80]{DejaVu Sans Mono}

%% ── Язык ────────────────────────────────────────────────────────────────────
\usepackage{polyglossia}
\setmainlanguage[babelshorthands=false]{russian}
\setotherlanguage{english}
\newfontfamily\cyrillicfont{Times New Roman}[Script=Cyrillic]
\newfontfamily\cyrillicfontsf{Times New Roman}[Script=Cyrillic]
\newfontfamily\cyrillicfonttt{DejaVu Sans Mono}[Script=Cyrillic,Scale=0.80]

%% ── Поля по ГОСТ 2.105-95 ────────────────────────────────────────────────────
\usepackage{geometry}
\geometry{left=30mm, right=15mm, top=20mm, bottom=20mm,
          headsep=6mm, footskip=12mm}

%% ── Межстрочный интервал ─────────────────────────────────────────────────────
\usepackage{setspace}
\onehalfspacing

%% ── Абзацный отступ ──────────────────────────────────────────────────────────
\usepackage{indentfirst}
\setlength{\parindent}{1.25cm}

%% ── Математика ───────────────────────────────────────────────────────────────
\usepackage{amsmath}
\usepackage{amssymb}

%% ── Таблицы ──────────────────────────────────────────────────────────────────
\usepackage{longtable}
\usepackage{booktabs}
\usepackage{tabularx}
\usepackage{array}
\usepackage{multirow}
\usepackage{pdflscape}
\newcolumntype{L}[1]{>{\raggedright\arraybackslash}p{#1}}
\newcolumntype{C}[1]{>{\centering\arraybackslash}p{#1}}
\newcolumntype{R}[1]{>{\raggedleft\arraybackslash}p{#1}}

%% ── Подписи к таблицам (ГОСТ: «Таблица N — Название») ───────────────────────
\usepackage{caption}
\DeclareCaptionLabelSeparator{gost}{ — }
\captionsetup[table]{
  name            = Таблица,
  labelsep        = gost,
  position        = top,
  justification   = raggedright,
  singlelinecheck = false,
  font            = normalsize,
  skip            = 4pt
}
\captionsetup[longtable]{
  name            = Таблица,
  labelsep        = gost,
  justification   = raggedright,
  singlelinecheck = false,
  font            = normalsize,
  skip            = 4pt
}

%% ── Листинги (каталоги, команды) ─────────────────────────────────────────────
\usepackage{listings}
\lstset{
  basicstyle   = \small\ttfamily,
  breaklines   = true,
  keepspaces   = true,
  columns      = flexible,
  frame        = single,
  framesep     = 4pt,
  xleftmargin  = 8pt,
  xrightmargin = 8pt,
  aboveskip    = 8pt,
  belowskip    = 8pt
}

%% ── Заголовки разделов ───────────────────────────────────────────────────────
\usepackage{titlesec}
\titleformat{\section}
  {\bfseries\normalsize}{\thesection\quad}{0pt}{}
\titlespacing*{\section}{0pt}{20pt}{10pt}

\titleformat{\subsection}
  {\bfseries\normalsize}{\thesubsection\quad}{0pt}{}
\titlespacing*{\subsection}{0pt}{14pt}{6pt}

\titleformat{\subsubsection}
  {\bfseries\normalsize}{\thesubsubsection\quad}{0pt}{}
\titlespacing*{\subsubsection}{0pt}{10pt}{4pt}

%% ── Оглавление ───────────────────────────────────────────────────────────────
\usepackage{tocloft}
\renewcommand{\cfttoctitlefont}{\hfill\bfseries\normalsize}
\renewcommand{\cftaftertoctitle}{\hfill}
\setlength{\cftbeforesecskip}{4pt}
\renewcommand{\cftsecfont}{\bfseries}
\renewcommand{\cftsecpagefont}{\bfseries}
\renewcommand{\cftdotsep}{1}
\setlength{\cftsecindent}{0pt}
\setlength{\cftsubsecindent}{1.5em}
\setlength{\cftsubsubsecindent}{3.5em}

%% ── Колонтитулы ──────────────────────────────────────────────────────────────
\usepackage{fancyhdr}
\pagestyle{fancy}
\fancyhf{}
\fancyhead[L]{\footnotesize Приложение ВКР — Описание набора данных и моделей}
\fancyfoot[C]{\thepage}
\renewcommand{\headrulewidth}{0.4pt}
\fancypagestyle{plain}{\fancyhf{}\fancyfoot[C]{\thepage}\renewcommand{\headrulewidth}{0pt}}

%% ── Гиперссылки ──────────────────────────────────────────────────────────────
\usepackage[colorlinks=false, hidelinks, unicode=true]{hyperref}

%% ── Вспомогательные команды ──────────────────────────────────────────────────
\newcommand{\tbl}[1]{\texttt{#1}}  % идентификаторы в тексте
\newcommand{\env}[1]{\texttt{#1}}  % переменные окружения

% ============================================================
\begin{document}

% ═══════════════════════════════════════════════════════════════════
% ТИТУЛЬНЫЙ ЛИСТ
% ═══════════════════════════════════════════════════════════════════
\begin{titlepage}
\thispagestyle{empty}
\begin{center}
  {\bfseries МИНИСТЕРСТВО НАУКИ И ВЫСШЕГО ОБРАЗОВАНИЯ\\
  РОССИЙСКОЙ ФЕДЕРАЦИИ}\vspace{4pt}

  {\small Федеральное государственное автономное образовательное
  учреждение высшего образования}

  \vspace{4pt}
  {\bfseries «САНКТ-ПЕТЕРБУРГСКИЙ ГОСУДАРСТВЕННЫЙ\\
  УНИВЕРСИТЕТ АЭРОКОСМИЧЕСКОГО ПРИБОРОСТРОЕНИЯ»}

  \vspace{4pt}(ГУАП)
\end{center}

\vspace{6pt}
\noindent\rule{\textwidth}{0.4pt}
\vspace{16pt}

\begin{center}
  {\bfseries AgroVision Production}

  \vspace{28pt}
  {\Large\bfseries ПРИЛОЖЕНИЕ ВКР}

  \vspace{10pt}
  {\bfseries ВЕБ-ПРИЛОЖЕНИЕ ИНФОРМАЦИОННОЙ ПОДДЕРЖКИ\\
  СЕЛЬСКОХОЗЯЙСТВЕННОЙ ДЕЯТЕЛЬНОСТИ\\
  НА ОСНОВЕ АГРОМЕТЕОРОЛОГИЧЕСКИХ ДАННЫХ}

  \vspace{20pt}
  {\Large\bfseries ОПИСАНИЕ НАБОРА ДАННЫХ И МОДЕЛЕЙ}

  \vspace{10pt}
  ПВКР.ОД.001-2026
\end{center}

\vspace{28pt}
\begin{tabular}{@{}lll@{}}
  \textbf{Разработал:} & техник-программист   & А.\,А.\,Майоров \\[6pt]
  \textbf{Проверил:}   & научный руководитель & П.\,А.\,Охтилев \\
\end{tabular}

\vfill
\begin{center}
  Санкт-Петербург — 2026
\end{center}
\end{titlepage}

% ═══════════════════════════════════════════════════════════════════
% ОГЛАВЛЕНИЕ
% ═══════════════════════════════════════════════════════════════════
\newpage
\tableofcontents
\newpage

% ═══════════════════════════════════════════════════════════════════
% АННОТАЦИЯ
% ═══════════════════════════════════════════════════════════════════
\section*{Аннотация}
\addcontentsline{toc}{section}{Аннотация}

Настоящий документ является описанием набора данных и программных моделей,
применяемых в программном комплексе \textbf{AGROSIGNAL} — веб-приложении
информационной поддержки сельскохозяйственной деятельности на основе
агрометеорологических данных.

Документ разработан в соответствии с требованиями ГОСТ~19.402-78
«Единая система программной документации. Описание программы»
и ГОСТ~2.105-95 «Единая система конструкторской документации.
Общие требования к текстовым документам».

Документ содержит техническое описание структуры обучающих данных, методологию
формирования признакового пространства, архитектуру нейросетевых моделей,
математические формулировки функций потерь, параметры инференса и методологию
регионального дообучения. Документ предназначен для разработчиков, осуществляющих
сопровождение и развитие системы, а также для научного руководства в рамках
выпускной квалификационной работы.

\textbf{Обозначение документа:} ПВКР.ОД.001-2026

% ═══════════════════════════════════════════════════════════════════
% НОРМАТИВНЫЕ ССЫЛКИ
% ═══════════════════════════════════════════════════════════════════
\section*{Нормативные ссылки}
\addcontentsline{toc}{section}{Нормативные ссылки}

При разработке настоящего документа использованы следующие нормативные документы:

\begin{enumerate}
  \item ГОСТ~19.402-78 «Единая система программной документации.
        Описание программы».
  \item ГОСТ~19.101-77 «Единая система программной документации.
        Виды программ и программных документов».
  \item ГОСТ~2.105-95 «Единая система конструкторской документации.
        Общие требования к текстовым документам».
  \item ГОСТ~34.201-89 «Информационная технология. Комплекс стандартов
        на автоматизированные системы. Виды, комплектность и обозначения
        документов при создании автоматизированных систем».
  \item РД~50-34.698-90 «Автоматизированные системы.
        Требования к содержанию документов».
\end{enumerate}

% ═══════════════════════════════════════════════════════════════════
% ТЕРМИНЫ И СОКРАЩЕНИЯ
% ═══════════════════════════════════════════════════════════════════
\section*{Термины и сокращения}
\addcontentsline{toc}{section}{Термины и сокращения}

\begin{longtable}{@{}L{30mm}L{120mm}@{}}
\toprule
\textbf{Термин} & \textbf{Определение} \\
\midrule
\endfirsthead
\multicolumn{2}{r}{\small(продолжение таблицы)}\\
\toprule
\textbf{Термин} & \textbf{Определение} \\
\midrule
\endhead
\bottomrule
\endfoot
\textbf{АОИ}       & Область интереса (Area of Interest) — географический регион,
                     в пределах которого выполняется обработка спутниковых данных. \\[4pt]
\textbf{BSI}       & Индекс голого почвенного покрова (Bare Soil Index) —
                     спектральный индекс, рассчитываемый по каналам Sentinel-2. \\[4pt]
\textbf{ГРД}       & Геопространственные растровые данные — данные, хранящие значения
                     пикселей в привязке к географическим координатам. \\[4pt]
\textbf{NDMI}      & Нормализованный индекс влажности
                     (Normalized Difference Moisture Index). \\[4pt]
\textbf{NDVI}      & Нормализованный разностный вегетационный индекс
                     (Normalized Difference Vegetation Index). \\[4pt]
\textbf{NDWI}      & Нормализованный разностный водный индекс
                     (Normalized Difference Water Index). \\[4pt]
\textbf{MNDWI}     & Модифицированный нормализованный разностный водный индекс
                     (Modified NDWI). \\[4pt]
\textbf{NPZ}       & Формат хранения массивов NumPy
                     (NumPy Zipped Archive). \\[4pt]
\textbf{ONNX}      & Открытый формат нейронных сетей
                     (Open Neural Network Exchange). \\[4pt]
\textbf{ПО}        & Программное обеспечение. \\[4pt]
\textbf{SAM}       & Segment Anything Model — нейросетевая модель универсальной
                     сегментации объектов (Meta AI). \\[4pt]
\textbf{SCL}       & Scene Classification Layer — растр классификации сцены
                     Sentinel-2 (классы 1--11). \\[4pt]
\textbf{Sentinel-2}& Спутниковая миссия ЕКА для наблюдения Земли с пространственным
                     разрешением 10--60~м. \\[4pt]
\textbf{TIF\,/\,GeoTIFF} & Формат географически привязанных растровых изображений. \\[4pt]
\textbf{UNet}      & Архитектура полностью свёрточной нейронной сети типа
                     «энкодер–декодер» с пропускными соединениями,
                     предназначенная для семантической сегментации. \\
\end{longtable}

% ═══════════════════════════════════════════════════════════════════
\section{Общие сведения}
% ═══════════════════════════════════════════════════════════════════

\subsection{Назначение набора данных}

Набор данных программного комплекса AGROSIGNAL предназначен для обучения
и верификации моделей автоматического выделения границ сельскохозяйственных
полей по многоспектральным снимкам спутника Sentinel-2. В состав набора входят
предобработанные тайлы спутниковых данных, метки полей, обученные нейросетевые
модели и вспомогательные конфигурационные файлы.

Система реализует следующие функции обработки данных:
\begin{itemize}
  \item автоматическое обнаружение и оцифровку границ сельскохозяйственных угодий
        по мультиспектральным снимкам Sentinel-2;
  \item постфильтрацию сегментированных объектов с применением
        ансамблевого классификатора;
  \item региональное дообучение основной модели сегментации на данных
        конкретных субъектов Российской Федерации.
\end{itemize}

\subsection{Сводные характеристики}

В таблице~\ref{tab:summary} приведены ключевые количественные характеристики
набора данных и моделей.

\begin{table}[h!]
\caption{Сводные характеристики программного комплекса}
\label{tab:summary}
\begin{tabularx}{\textwidth}{L{105mm}L{50mm}}
\toprule
\textbf{Характеристика} & \textbf{Значение} \\
\midrule
Метрика IoU Extent (выборка валидации)           & 0,534 \\
Метрика F1 Boundary (выборка валидации)          & 0,581 \\
Расстояние Хаусдорфа 95-й перцентиль            & 14,27~пикс. \\
Отклонение ONNX от эталона (max delta)           & менее $5{,}1\!\cdot\!10^{-8}$ \\
Число обучаемых параметров модели                & 7~766~851 \\
Размерность входного признакового пространства   & 16 каналов (профиль v2\_16ch) \\
Число тайлов глобального датасета                & 64 \\
Число региональных тайлов                        & 22 (12~— Краснодарский край, 10~— Пермский край) \\
Число обучающих патчей                           & 7~774 \\
Число патчей выборки валидации                   & 1~690 \\
Число патчей тестовой выборки                    & 1~352 \\
\bottomrule
\end{tabularx}
\end{table}

% ═══════════════════════════════════════════════════════════════════
\section{Структура хранения данных}
% ═══════════════════════════════════════════════════════════════════

\subsection{Локальное файловое хранилище}

Данные программного комплекса размещаются в файловой системе сервера
в соответствии со следующей структурой каталогов:

\begin{lstlisting}
backend/
├── debug/
│   ├── runs/
│   │   ├── real_tiles/                  # обработанные NPZ-тайлы Sentinel-2
│   │   │   ├── tile_*.npz               # глобальный датасет (64 тайла)
│   │   │   ├── krd_*.npz                # Краснодарский край (12 тайлов)
│   │   │   └── permkrai_*.npz           # Пермский край (10 тайлов)
│   │   └── real_tiles_labels_weak/      # растровые метки полей (GeoTIFF)
│   │       ├── tile_*_label.tif         # метки глобального датасета (64)
│   │       ├── krd_*_label.tif          # метки Краснодарского края (12)
│   │       └── permkrai_*_label.tif     # метки Пермского края (10)
│   └── cache/
│       └── sentinel_scenes/             # кэш исходных сцен Sentinel Hub
└── models/
    ├── boundary_unet_v3_cpu.pth         # чекпойнт модели PyTorch
    ├── boundary_unet_v3_cpu.onnx        # экспорт модели в формате ONNX
    ├── boundary_unet_v3_cpu.norm.json   # статистики нормализации признаков
    ├── object_classifier.pkl            # сериализованная модель классификатора
    ├── sam_vit_h_4b8939.pth             # модель SAM ViT-H
    └── regional_training_report_*.json  # отчёты регионального обучения
\end{lstlisting}

\subsection{Облачное хранилище Google Drive}

Резервное и дистрибутивное хранение данных осуществляется посредством сервиса
Google Drive с использованием инструмента rclone. Идентификатор корневой
директории: \tbl{1jtUAN8KRcJb73dxTME-Gi1d4TvGF3P2z}.

\begin{lstlisting}
gdrive:/
├── models/                              # модели (~3 ГБ)
│   ├── boundary_unet_v3_cpu.pth
│   ├── boundary_unet_v3_cpu.onnx
│   ├── boundary_unet_v3_cpu.norm.json
│   └── object_classifier.pkl
├── training/
│   ├── real_tiles/           # глобальный датасет (64 NPZ-файла, ~2,7 ГБ)
│   └── real_tiles_labels_weak/  # метки глобального датасета (64 TIF, ~800 МБ)
├── cache/
│   └── sentinel_scenes/      # кэш исходных сцен (~26 ГБ, не обязателен)
└── artifacts/
    └── agrosignal_dataset_20260316.tar  # резервный архив (~2,7 ГБ)
\end{lstlisting}

Региональные тайлы (22 единицы) в облачном хранилище не размещаются
и при необходимости воспроизводятся посредством повторного запуска пайплайна
\tbl{train\_regional\_supplement.py}.

\subsection{Управление данными облачного хранилища}

Для управления данными между локальным хранилищем и Google Drive используется
скрипт \tbl{scripts/gdrive\_data.py}. Доступные операции приведены
в таблице~\ref{tab:gdrive_commands}.

\begin{table}[h!]
\caption{Команды управления данными}
\label{tab:gdrive_commands}
\begin{tabularx}{\textwidth}{L{78mm}X}
\toprule
\textbf{Команда} & \textbf{Описание} \\
\midrule
\tbl{python scripts/gdrive\_data.py status}              & Отображение состояния всех наборов данных \\
\tbl{python scripts/gdrive\_data.py pull models}         & Загрузка моделей из облачного хранилища \\
\tbl{python scripts/gdrive\_data.py pull real\_tiles}    & Загрузка глобального датасета \\
\tbl{python scripts/gdrive\_data.py pull real\_tiles\_labels\_weak} & Загрузка меток глобального датасета \\
\tbl{python scripts/gdrive\_data.py push models}         & Выгрузка обновлённых моделей в облачное хранилище \\
\tbl{python scripts/gdrive\_data.py pull {-}{-}dry-run}      & Предварительный просмотр операции без выполнения \\
\bottomrule
\end{tabularx}
\end{table}

\begin{table}[h!]
\caption{Ключи наборов данных}
\label{tab:data_keys}
\begin{tabularx}{\textwidth}{L{48mm}L{20mm}X}
\toprule
\textbf{Ключ} & \textbf{Объём} & \textbf{Описание} \\
\midrule
\tbl{real\_tiles}               & ${\sim}2{,}7$~ГБ & 64 NPZ-тайла Sentinel-2 (глобальный датасет) \\
\tbl{real\_tiles\_labels\_weak} & ${\sim}800$~МБ   & 64 GeoTIFF-файла слабых меток полей \\
\tbl{sentinel\_cache}           & ${\sim}26$~ГБ    & Кэш исходных сцен (не обязателен для инференса) \\
\tbl{models}                    & ${\sim}3$~ГБ     & Все обученные модели и вспомогательные файлы \\
\tbl{dataset\_tar}              & ${\sim}2{,}7$~ГБ & Резервный архив обучающего датасета \\
\bottomrule
\end{tabularx}
\end{table}

% ═══════════════════════════════════════════════════════════════════
\section{Спутниковые данные Sentinel-2}
% ═══════════════════════════════════════════════════════════════════

\subsection{Характеристики тайлов}

Тайл представляет собой предобработанный фрагмент многоспектрального снимка
Sentinel-2, хранимый в формате NPZ. Каждый тайл охватывает область
$512 \times 512$~пикселей при пространственном разрешении 10~м/пиксель,
что соответствует площади приблизительно 26,2~км². Для получения данных
используются спектральные каналы B2~(490~нм), B3~(560~нм), B4~(665~нм),
B8~(842~нм), B11~(1610~нм), B12~(2190~нм), а также слой классификации
сцены SCL.

\subsubsection{Метод формирования стека признаков}

Признаковый стек формируется на основе многовременной последовательности
снимков Sentinel-2 за 12~месяцев, предшествующих дате создания тайла.
Процедура включает следующие шаги:

\begin{enumerate}
  \item \textbf{Отбор сцен.} Из архива Sentinel Hub отбираются сцены,
        для которых доля облачного покрова по каналу SCL не превышает порогового
        значения, определяемого региональным профилем обработки.

  \item \textbf{Маскирование облачности.} Из каждой сцены исключаются пиксели
        с кодами SCL, соответствующими облакам (код~8), теням от облаков (код~3),
        снегу (код~11) и насыщению сенсора (код~1). Валидными считаются пиксели
        с кодами 4 (растительность), 5 (голый грунт), 6 (водная поверхность),
        7 (неклассифицированные).

  \item \textbf{Вычисление спектральных индексов.} Для каждой валидной сцены
        рассчитываются NDVI, NDWI, NDMI, MNDWI, BSI.

  \item \textbf{Временн\'{а}я агрегация.} По временной оси вычисляются
        статистические агрегаты: максимум, среднее, стандартное отклонение,
        медиана и их комбинации. Результат — 16-канальный растр
        размером $512 \times 512$~пикселей.

  \item \textbf{Формирование признака рёбер.} Канал \tbl{edge\_composite}
        вычисляется применением оператора Собела к нормированному цветному
        composite (B4, B3, B2) с последующим максимумом по каналам.
\end{enumerate}

Итоговый стек нормализуется per-channel z-нормализацией
с параметрами из файла \tbl{boundary\_unet\_v3\_cpu.norm.json}
(см.~раздел~\ref{sec:norm}) перед подачей в модель.

\subsection{Структура NPZ-файла}

Каждый файл формата \tbl{.npz} содержит набор именованных массивов,
перечисленных в таблице~\ref{tab:npz}.

\begin{table}[h!]
\caption{Состав NPZ-файла тайла Sentinel-2}
\label{tab:npz}
\begin{tabularx}{\textwidth}{L{30mm}L{26mm}L{14mm}X}
\toprule
\textbf{Ключ} & \textbf{Форма} & \textbf{Тип} & \textbf{Содержание} \\
\midrule
\tbl{features}    & (16, 512, 512) & float32 & Стек вычисленных признаков (16 каналов) \\
\tbl{B2}          & (512, 512)     & float32 & Отражение в синем канале \\
\tbl{B3}          & (512, 512)     & float32 & Отражение в зелёном канале \\
\tbl{B4}          & (512, 512)     & float32 & Отражение в красном канале \\
\tbl{B8}          & (512, 512)     & float32 & Отражение в ближнем ИК-диапазоне \\
\tbl{B11}         & (512, 512)     & float32 & Отражение в SWIR-1 диапазоне \\
\tbl{B12}         & (512, 512)     & float32 & Отражение в SWIR-2 диапазоне \\
\tbl{SCL}         & (512, 512)     & uint8   & Классификация сцены (классы 1--11) \\
\tbl{bbox}        & (4,)           & float64 & Прямоугольник [lon\_min, lat\_min, lon\_max, lat\_max] \\
\tbl{crs}         & скаляр         & str     & Система координат (EPSG:4326) \\
\tbl{date\_from}  & скаляр         & str     & Начало временного окна наблюдения \\
\tbl{date\_to}    & скаляр         & str     & Конец временного окна наблюдения \\
\tbl{tile\_id}    & скаляр         & str     & Уникальный идентификатор тайла \\
\tbl{ndvi}        & (512, 512)     & float32 & $\text{NDVI} = (B8 - B4)/(B8 + B4)$ \\
\tbl{ndwi}        & (512, 512)     & float32 & $\text{NDWI} = (B3 - B8)/(B3 + B8)$ \\
\tbl{ndmi}        & (512, 512)     & float32 & $\text{NDMI} = (B8 - B11)/(B8 + B11)$ \\
\tbl{valid\_mask} & (512, 512)     & bool    & Маска валидных (безоблачных) пикселей \\
\tbl{cloud\_pct}  & скаляр         & float32 & Доля облачного покрова сцены, \% \\
\tbl{source}      & скаляр         & str     & Источник данных (\tbl{sentinel\_hub} / \tbl{synthetic}) \\
\tbl{n\_dates}    & скаляр         & int32   & Число используемых временных срезов \\
\tbl{profile\_key}& скаляр         & str     & Ключ регионального профиля обработки \\
\bottomrule
\end{tabularx}
\end{table}

\subsection{Признаковое пространство (профиль v2\_16ch)}
\label{sec:features}

Признаки вычисляются на основе временн\'{о}го ряда снимков Sentinel-2
за 12~месяцев. В таблице~\ref{tab:features} приведены наименования,
физический смысл и статистики нормализации всех 16~каналов.

Признаковое пространство разработано с учётом следующих принципов:
\begin{itemize}
  \item \textbf{Фенологические признаки} (\tbl{max\_ndvi}, \tbl{mean\_ndvi},
        \tbl{ndvi\_std}) отражают сезонную динамику растительности
        и позволяют отличить регулярно обрабатываемые поля от прочих
        землепользований. Поля характеризуются повышенными значениями
        \tbl{max\_ndvi} в вегетационный период и значительной сезонной
        изменчивостью (\tbl{ndvi\_std}).

  \item \textbf{Водные и влажностные индексы} (\tbl{ndwi\_mean},
        \tbl{mndwi\_max}, \tbl{ndmi\_mean}, \tbl{ndwi\_median}) обеспечивают
        разделение рисовых чеков, орошаемых угодий и богарных полей.
        Высокие значения \tbl{mndwi\_max} характерны для постоянных
        водных объектов.

  \item \textbf{Почвенные индексы} (\tbl{bsi\_mean}) выделяют периоды
        парования и ранние фенологические фазы, когда поле представляет
        собой голый грунт.

  \item \textbf{Текстурные признаки} (\tbl{ndvi\_entropy},
        \tbl{edge\_composite}) несут информацию о пространственной
        однородности. Сельскохозяйственные поля, как правило, обладают
        высокой внутренней однородностью и резкими границами,
        что выражается в повышенных значениях \tbl{edge\_composite}
        вблизи контуров.

  \item \textbf{Яркостные признаки} (\tbl{rgb\_r}, \tbl{rgb\_g},
        \tbl{rgb\_b}, \tbl{green\_median}, \tbl{swir\_median}) кодируют
        среднее спектральное отражение, полезное для различения типов
        культур и почвенного покрова.
\end{itemize}

\begin{landscape}
\begin{longtable}{@{}C{6mm}L{28mm}L{44mm}L{48mm}R{14mm}R{16mm}@{}}
\caption{Состав признакового пространства v2\_16ch\label{tab:features}}\\
\toprule
\textbf{№} & \textbf{Наименование} & \textbf{Физический смысл} &
\textbf{Математическое определение} & \textbf{Среднее} & \textbf{Стд. откл.} \\
\midrule
\endfirsthead
\multicolumn{6}{r}{\small(продолжение таблицы~\ref{tab:features})}\\
\toprule
\textbf{№} & \textbf{Наименование} & \textbf{Физический смысл} &
\textbf{Математическое определение} & \textbf{Среднее} & \textbf{Стд. откл.} \\
\midrule
\endhead
\bottomrule
\endfoot
0 & \tbl{edge\_composite}    & Карта рёбер — градиент яркости, выделяющий границы полей
  & Оператор Собела по RGB-composite, max по каналам                        & 0,2537 & 0,2944 \\[3pt]
1 & \tbl{max\_ndvi}          & Максимальный NDVI за вегетационный сезон
  & $\max_{t}(\text{NDVI}_t)$                                               & 0,7061 & 0,2273 \\[3pt]
2 & \tbl{mean\_ndvi}         & Среднегодовой NDVI — интегральная продуктивность
  & $\overline{\text{NDVI}}$                                                 & 0,4591 & 0,2140 \\[3pt]
3 & \tbl{ndvi\_std}          & Стандартное отклонение NDVI — сезонная динамика
  & $\sigma(\text{NDVI}_t)$                                                  & 0,1753 & 0,0757 \\[3pt]
4 & \tbl{ndwi\_mean}         & Среднегодовой NDWI — влажность растительности
  & $\overline{\text{NDWI}}$, $\text{NDWI}=(B3-B8)/(B3+B8)$                 & $-0,4721$ & 0,2038 \\[3pt]
5 & \tbl{bsi\_mean}          & Средний индекс голого почвенного покрова
  & $\overline{\text{BSI}}$, $\text{BSI}=((B11{+}B4)-(B8{+}B2))/((B11{+}B4)+(B8{+}B2))$
  & 0,0006 & 0,1389 \\[3pt]
6 & \tbl{scl\_valid\_fraction}& Доля пикселей без облаков и теней по SCL
  & $|\{p : \text{SCL}_p \in \{4,5,6,7\}\}| / N$                           & 0,6693 & 0,1180 \\[3pt]
7 & \tbl{rgb\_r}             & Медианное красное отражение (B4)
  & $\text{median}_t(B4_t)$                                                  & 0,2080 & 0,0716 \\[3pt]
8 & \tbl{rgb\_g}             & Медианное зелёное отражение (B3)
  & $\text{median}_t(B3_t)$                                                  & 0,0652 & 0,0337 \\[3pt]
9 & \tbl{rgb\_b}             & Медианное синее отражение (B2)
  & $\text{median}_t(B2_t)$                                                  & 0,0457 & 0,0232 \\[3pt]
10& \tbl{ndvi\_entropy}      & Энтропия распределения NDVI — текстурная однородность
  & $H(\text{NDVI})$ по гистограмме                                          & 1,7698 & 0,5598 \\[3pt]
11& \tbl{mndwi\_max}         & Максимальный MNDWI — наличие водных объектов
  & $\max_t(\text{MNDWI}_t)$, $\text{MNDWI}=(B3-B11)/(B3+B11)$             & $-0,2850$ & 0,3092 \\[3pt]
12& \tbl{ndmi\_mean}         & Средний NDMI — влажность растительного покрова
  & $\overline{\text{NDMI}}$, $\text{NDMI}=(B8-B11)/(B8+B11)$              & 0,0372 & 0,1606 \\[3pt]
13& \tbl{ndwi\_median}       & Медианный NDWI — устойчивая характеристика влажности
  & $\text{median}_t(\text{NDWI}_t)$                                        & $-0,4731$ & 0,2156 \\[3pt]
14& \tbl{green\_median}      & Медианное зелёное отражение B3 по чистым пикселям
  & $\text{median}_{t,\text{valid}}(B3_t)$                                   & 0,0653 & 0,0262 \\[3pt]
15& \tbl{swir\_median}       & Медианное отражение в SWIR-1 диапазоне (B11)
  & $\text{median}_t(B11_t)$                                                 & 0,1920 & 0,0634 \\
\end{longtable}
\end{landscape}

\subsection{Глобальный датасет}

\subsubsection{Источники эталонных данных}

Формирование меток глобального датасета выполнено с использованием четырёх
открытых источников эталонной разметки, указанных в таблице~\ref{tab:gt_sources}.

\begin{table}[h!]
\caption{Источники эталонных данных глобального датасета}
\label{tab:gt_sources}
\begin{tabularx}{\textwidth}{L{44mm}L{60mm}X}
\toprule
\textbf{Источник} & \textbf{Описание} & \textbf{Охватываемый регион} \\
\midrule
AI4Boundaries           & Официальные контуры сельскохозяйственных угодий ЕС   & Западная и Центральная Европа \\
EuroCropsV2             & Мультивременные метки сельскохозяйственных культур   & Европа \\
Fields of The World (FoTW) & Разнообразные контуры полей мирового охвата       & Глобально \\
Russia/Ukraine dataset  & Контуры полей на территории России и Украины         & Восточная Европа \\
\bottomrule
\end{tabularx}
\end{table}

\subsubsection{Состав выборок}

Разбиение датасета на обучающую, валидационную и тестовую выборки произведено
на уровне тайлов с соблюдением географической диверсификации (6~кластеров
по климатическим зонам: Центральная Европа, Западная Европа, Северная Европа,
Балканы, Россия/Украина, Африка/Южная Азия). Географическое разнесение
выборок исключает утечку данных между обучением и верификацией при наличии
пространственной автокорреляции.

Размеры патчей составляют $256 \times 256$~пикселей, шаг экстракции —
128~пикселей (перекрытие 50\,\%).

\begin{table}[h!]
\caption{Состав обучающего датасета}
\label{tab:splits}
\begin{tabularx}{\textwidth}{Xrr}
\toprule
\textbf{Выборка} & \textbf{Число тайлов} & \textbf{Число патчей} \\
\midrule
Обучающая    & 46 & 7~774 \\
Валидационная & 10 & 1~690 \\
Тестовая      & 8  & 1~352 \\
\midrule
\textbf{Итого}& \textbf{64} & \textbf{10~816} \\
\bottomrule
\end{tabularx}
\end{table}

\subsection{Региональные тайлы — Краснодарский край}

Региональный датасет для Краснодарского края (код ОКТМО~23) сформирован
с целью дообучения модели на характерных сельскохозяйственных условиях
юга России: рисоводство, озимая пшеница, подсолнечник, кукуруза.
Состав регионального датасета представлен в таблице~\ref{tab:krd_tiles}.

\begin{landscape}
\begin{table}[h!]
\caption{Тайлы Краснодарского края}
\label{tab:krd_tiles}
\begin{tabularx}{\linewidth}{L{40mm}R{16mm}R{16mm}R{16mm}R{16mm}X}
\toprule
\textbf{Идентификатор} & \textbf{lon\_min} & \textbf{lat\_min} &
\textbf{lon\_max} & \textbf{lat\_max} & \textbf{Описание зоны} \\
\midrule
\tbl{krd\_rice\_e}         & 39,0 & 45,1 & 39,5 & 45,5 & Рисовые чеки (восточная зона) \\
\tbl{krd\_rice\_w}         & 38,3 & 45,2 & 38,8 & 45,6 & Рисовые чеки (западная зона) \\
\tbl{krd\_wheat\_n}        & 38,5 & 46,5 & 39,0 & 46,9 & Озимая пшеница (северная зона) \\
\tbl{krd\_wheat\_s}        & 38,7 & 44,9 & 39,2 & 45,3 & Озимая пшеница (южная зона) \\
\tbl{krd\_sun\_e}          & 39,5 & 45,8 & 40,0 & 46,2 & Подсолнечник (восточная зона) \\
\tbl{krd\_sun\_w}          & 38,0 & 46,0 & 38,5 & 46,4 & Подсолнечник (западная зона) \\
\tbl{krd\_corn\_n}         & 38,9 & 45,7 & 39,4 & 46,1 & Кукуруза (северная зона) \\
\tbl{krd\_corn\_s}         & 39,2 & 44,8 & 39,7 & 45,2 & Кукуруза (южная зона) \\
\tbl{krd\_mix\_krasnodar}  & 38,9 & 45,0 & 39,4 & 45,4 & Смешанные культуры, район Краснодара \\
\tbl{krd\_mix\_tikhoretsk} & 40,1 & 45,8 & 40,6 & 46,2 & Смешанные культуры, район Тихорецка \\
\tbl{krd\_mix\_kropotkin}  & 40,5 & 45,4 & 41,0 & 45,8 & Смешанные культуры, район Кропоткина \\
\tbl{krd\_mix\_temruk}     & 37,3 & 45,2 & 37,8 & 45,6 & Смешанные культуры, Таманский п-ов \\
\bottomrule
\end{tabularx}
\end{table}
\end{landscape}

\subsection{Региональные тайлы — Пермский край}

Региональный датасет для Пермского края (код ОКТМО~59) сформирован для условий
таёжной зоны Предуралья: мелкоконтурные поля, зерновые культуры, кормовые травы.
Состав регионального датасета представлен в таблице~\ref{tab:perm_tiles}.

\begin{landscape}
\begin{table}[h!]
\caption{Тайлы Пермского края}
\label{tab:perm_tiles}
\begin{tabularx}{\linewidth}{L{44mm}R{16mm}R{16mm}R{16mm}R{16mm}X}
\toprule
\textbf{Идентификатор} & \textbf{lon\_min} & \textbf{lat\_min} &
\textbf{lon\_max} & \textbf{lat\_max} & \textbf{Административный район} \\
\midrule
\tbl{permkrai\_barda}  & 56,2 & 56,9 & 56,7 & 57,3 & Бардымский район \\
\tbl{permkrai\_chaikov}& 54,0 & 56,7 & 54,5 & 57,1 & Чайковский район \\
\tbl{permkrai\_kungur} & 56,9 & 57,3 & 57,4 & 57,7 & Кунгурский район \\
\tbl{permkrai\_ocher}  & 54,7 & 57,8 & 55,2 & 58,2 & Очёрский район \\
\tbl{permkrai\_ordin}  & 55,5 & 57,0 & 56,0 & 57,4 & Ординский район \\
\tbl{permkrai\_osa}    & 55,4 & 57,2 & 55,9 & 57,6 & Осинский район \\
\tbl{permkrai\_veresh} & 54,6 & 58,0 & 55,1 & 58,4 & Верещагинский район \\
\tbl{permkrai\_perm\_s}& 56,1 & 57,8 & 56,6 & 58,2 & Южная пригородная зона г.\,Перми \\
\tbl{permkrai\_dobr}   & 55,8 & 57,5 & 56,3 & 57,9 & Добрянский район \\
\tbl{permkrai\_nytva}  & 55,3 & 58,0 & 55,8 & 58,4 & Нытвенский район \\
\bottomrule
\end{tabularx}
\end{table}
\end{landscape}

% ═══════════════════════════════════════════════════════════════════
\section{Метки полей}
% ═══════════════════════════════════════════════════════════════════

\subsection{Формат хранения}

Метки полей хранятся в формате GeoTIFF (один растровый канал, тип данных
uint8) в координатной привязке EPSG:4326, согласованной с соответствующим
NPZ-тайлом. Значения пикселей кодируются следующим образом:
значение \tbl{0}~— фон (отсутствие сельскохозяйственного поля);
значение \tbl{1}~— сельскохозяйственное поле (слабая метка).

\subsection{Метод формирования слабых меток}

Формирование меток выполняется скриптом
\tbl{backend/training/generate\_weak\_labels\_real\_tiles.py}
посредством следующей последовательности операций:

\begin{enumerate}
  \item Запрос к Overpass API (OpenStreetMap) для получения полигонов
        с тегами \tbl{landuse=farmland}, \tbl{landuse=meadow},
        \tbl{landuse=orchard} в пределах ограничивающего прямоугольника тайла.
  \item Растеризация полученных полигонов в пиксельную сетку тайла (EPSG:4326).
  \item Морфологическая постобработка: дилатация и эрозия для устранения
        краевых артефактов.
\end{enumerate}

\textbf{Влияние качества слабых меток на процесс обучения.}
Данные OpenStreetMap не обеспечивают пространственную точность, соответствующую
разрешению Sentinel-2~(10~м). Типичное плановое смещение полигонов составляет
5--15~м, что на уровне одного пикселя влечёт систематическую погрешность
при обучении канала Boundary. Для компенсации данного эффекта применяется
следующая совокупность мер:
\begin{itemize}
  \item повышенный вес $w_b$ функции потерь для канала Boundary,
        что снижает относительный вклад пограничных пикселей с неопределённой
        меткой в итоговый градиент;
  \item использование морфологического градиента при формировании целевой
        карты $y_b$ (см.~раздел~\ref{sec:targets}), что уширяет область
        истинных граничных пикселей и делает метку менее чувствительной
        к плановому смещению;
  \item сильная аугментация (упругие деформации, случайные повороты)
        для обучения модели инвариантности к небольшим позиционным сдвигам.
\end{itemize}

Доля пикселей-полей в обучающей выборке составляет 15--30\,\%, что создаёт
умеренный дисбаланс классов. Он компенсируется слагаемым Dice Loss в функции
потерь канала Extent, которое по построению инвариантно к дисбалансу
(см.~раздел~\ref{sec:loss}).

\subsection{Характеристики покрытия меток по регионам}

Степень покрытия меток существенно варьируется в зависимости от региона
и определяется полнотой разметки в OpenStreetMap.
В таблице~\ref{tab:label_coverage} приведены типичные диапазоны покрытия.

\begin{table}[h!]
\caption{Типичное покрытие меток по регионам}
\label{tab:label_coverage}
\begin{tabularx}{\textwidth}{L{50mm}C{30mm}X}
\toprule
\textbf{Регион} & \textbf{Покрытие, \%} & \textbf{Причина вариации} \\
\midrule
Краснодарский край & 15--65 & Высокая плотность и качество разметки OSM \\
Пермский край      & 1--20  & Мелкоконтурные поля, ограниченная разметка OSM \\
Глобальный датасет & 5--55  & Зависит от страны и источника GT \\
\bottomrule
\end{tabularx}
\end{table}

% ═══════════════════════════════════════════════════════════════════
\section{Описание моделей}
% ═══════════════════════════════════════════════════════════════════

\subsection{Модель сегментации BoundaryUNet v3}

\subsubsection{Архитектура}

Основная модель сегментации реализована на основе архитектуры UNet с энкодером
и декодером, связанными пропускными соединениями (skip-connections). Модель
принимает на вход тензор $\mathbb{R}^{16\times512\times512}$ и формирует
три выходных карты вероятностей:
\begin{itemize}
  \item \textbf{Канал~0~— Extent:} вероятность принадлежности пикселя
        к площади сельскохозяйственного поля;
  \item \textbf{Канал~1~— Boundary:} вероятность принадлежности пикселя
        к границе поля;
  \item \textbf{Канал~2~— Distance:} нормализованное расстояние до ближайшей
        границы поля.
\end{itemize}

\begin{lstlisting}
Энкодер:
  conv0:  16 -> 32 ch,  stride 1  (свёртка 3x3, BatchNorm, ReLU)
  conv1:  32 -> 64 ch,  stride 2  (MaxPool + двойная свёртка)
  conv2:  64 -> 128 ch, stride 2
  conv3: 128 -> 256 ch, stride 2

Бутылочное горлышко:
  conv4: 256 -> 512 ch  (Dropout p=0.30)

Декодер (с пропускными соединениями):
  up3: 512+256 -> 256 ch  (Dropout p=0.20)
  up2: 256+128 -> 128 ch  (Dropout p=0.10)
  up1: 128+64  ->  64 ch
  up0:  64+32  ->  32 ch

Выходной блок:
  out: 32 -> 3 ch  (функция активации Sigmoid)
\end{lstlisting}

Общее число обучаемых параметров составляет \textbf{7~766~851}.

\textbf{Функциональная роль архитектурных элементов.}
Нормализация BatchNorm в каждом блоке сглаживает распределение активаций
и стабилизирует обучение при высокой скорости. Пропускные соединения
передают признаки высокого пространственного разрешения из энкодера
в декодер, что критически важно для точной локализации границ полей
на уровне одного пикселя. Dropout в бутылочном горлышке (p\,=\,0,30)
снижает переобучение на наиболее абстрактных признаках, где рецептивное
поле модели максимально. Пониженные коэффициенты Dropout в декодере
(0,20 и 0,10) обеспечивают сохранение пространственной детализации.
Многозадачное обучение по трём каналам одновременно улучшает качество
каждого в отдельности за счёт общего представления в энкодере.

\subsubsection{Файлы модели}

\begin{table}[h!]
\caption{Файлы модели BoundaryUNet v3}
\label{tab:model_files}
\begin{tabularx}{\textwidth}{L{68mm}L{16mm}X}
\toprule
\textbf{Файл} & \textbf{Размер} & \textbf{Назначение} \\
\midrule
\tbl{boundary\_unet\_v3\_cpu.pth}       & ${\sim}30$~МБ & Чекпойнт модели в формате PyTorch \\
\tbl{boundary\_unet\_v3\_cpu.onnx}      & ${\sim}30$~МБ & Экспорт модели в формате ONNX (opset~18) \\
\tbl{boundary\_unet\_v3\_cpu.norm.json} & $<$1~МБ       & Статистики нормализации признаков \\
\bottomrule
\end{tabularx}
\end{table}

\subsubsection{Параметры нормализации входных данных}
\label{sec:norm}

Перед подачей на вход модели признаки нормализуются по формуле
z-нормализации:
\begin{equation}
  \hat{x}_c = \frac{x_c - \mu_c}{\sigma_c},
  \label{eq:znorm}
\end{equation}
где $x_c$ — значение $c$-го канала, $\mu_c$ и $\sigma_c$ — математическое
ожидание и стандартное отклонение канала, хранящиеся в файле
\tbl{boundary\_unet\_v3\_cpu.norm.json}. Параметры приведены
в таблице~\ref{tab:features} (раздел~\ref{sec:features}).

Дополнительные параметры конфигурации модели:
\begin{itemize}
  \item порог классификации по каналу Extent: $\tau = 0{,}42$;
  \item версия ONNX opset: 18;
  \item $p_{\text{drop,bottleneck}} = 0{,}30$;\quad
        $p_{\text{drop,up3}} = 0{,}20$;\quad $p_{\text{drop,up2}} = 0{,}10$.
\end{itemize}

\subsubsection{Формирование целевых карт обучающей выборки}
\label{sec:targets}

Из бинарной маски поля $y_e \in \{0,1\}^{H\times W}$, полученной
из слабых меток OSM, дополнительно формируются:

\begin{itemize}
  \item \textbf{Карта границ} $y_b$ — морфологический градиент маски:
        \begin{equation}
          y_b = \text{dil}_r(y_e) \oplus \text{ero}_r(y_e),
          \label{eq:boundary_map}
        \end{equation}
        где $\text{dil}_r$ и $\text{ero}_r$ — дилатация и эрозия
        дисковым структурным элементом радиуса $r$~(пикс.), $\oplus$~—
        операция симметрической разности.

  \item \textbf{Карта расстояний} $y_d$ — нормализованное евклидово
        расстояние от каждого пикселя площади поля до ближайшей границы:
        \begin{equation}
          y_d = \text{norm}\!\left(\text{EDT}(\mathbf{1} - y_b)\right) \odot y_e,
          \label{eq:distance_map}
        \end{equation}
        где $\text{EDT}$ — евклидово дистанционное преобразование,
        $\text{norm}(\cdot)$ — нормирование в диапазон $[0, 1]$,
        $\odot$ — поэлементное произведение.
\end{itemize}

\subsubsection{Функция потерь}
\label{sec:loss}

Суммарная функция потерь задана как взвешенная сумма трёх слагаемых:
\begin{equation}
  \mathcal{L} = w_e \cdot \mathcal{L}_{e} + w_b \cdot \mathcal{L}_{b}
              + w_d \cdot \mathcal{L}_{d},
  \label{eq:loss_total}
\end{equation}

где веса $w_e$, $w_b$, $w_d$ задаются региональным профилем
(см.~раздел~\ref{sec:profiles}).

\textbf{Слагаемое для канала Extent:}
\begin{equation}
  \mathcal{L}_{e} = \alpha\,\text{BCE}(\hat{p}_e, y_e)
                  + \beta\,\mathcal{L}_{\text{Dice}}(\hat{p}_e, y_e),
  \label{eq:loss_extent}
\end{equation}
\begin{equation}
  \mathcal{L}_{\text{Dice}} = 1 - \frac{2\sum_{i}\hat{p}_i y_i}
                                       {\sum_{i}\hat{p}_i + \sum_{i}y_i + \varepsilon},
  \label{eq:dice}
\end{equation}
где $\varepsilon = 10^{-6}$ — сглаживающая константа, предотвращающая
деление на нуль. Dice Loss инвариантно к дисбалансу классов, что обосновывает
его использование совместно с BCE.

\textbf{Слагаемые для каналов Boundary и Distance:}
\begin{equation}
  \mathcal{L}_{b} = \text{BCE}(\hat{p}_b,\, y_b),
  \label{eq:loss_boundary}
\end{equation}
\begin{equation}
  \mathcal{L}_{d} = \text{MSE}(\hat{p}_d,\, y_d).
  \label{eq:loss_distance}
\end{equation}

\subsubsection{Параметры обучения базовой модели}

Обучение базовой модели на глобальном датасете выполнено со следующими
параметрами:
\begin{itemize}
  \item оптимизатор: AdamW;
  \item планировщик скорости обучения: CosineAnnealingLR;
  \item аугментация: случайное горизонтальное и вертикальное отражение,
        повороты на произвольный угол, упругая деформация с параметром
        $\alpha = 120$, случайное кадрирование до размера $256\times256$;
  \item функция потерь~\eqref{eq:loss_total}: $w_e = 1{,}0$,
        $w_b = 2{,}0$, $w_d = 0{,}5$;
        $\alpha = 0{,}5$, $\beta = 0{,}5$ в~\eqref{eq:loss_extent}.
\end{itemize}

\subsection{Классификатор объектов}

\subsubsection{Описание}

Классификатор объектов предназначен для постфильтрации сегментов,
выделенных моделью BoundaryUNet. Классификатор принимает на вход вектор
из 19~признаков, описывающих геометрические и спектральные характеристики
объекта, и выдаёт оценку вероятности того, что объект является
сельскохозяйственным полем.

\textit{Тип модели:} \tbl{sklearn.ensemble.HistGradientBoostingClassifier}.
\textit{Файл:} \tbl{backend/models/object\_classifier.pkl}.
\textit{Порог классификации} (параметр \env{OBJECT\_MIN\_SCORE}):
$\tau_{\text{cls}} = 0{,}40$.

\subsubsection{Входные признаки классификатора}

\begin{table}[h!]
\caption{Входные признаки классификатора объектов}
\label{tab:classifier_features}
\begin{tabularx}{\textwidth}{C{8mm}L{42mm}X}
\toprule
\textbf{№} & \textbf{Признак} & \textbf{Описание} \\
\midrule
0  & \tbl{area\_ha}            & Площадь объекта, га \\
1  & \tbl{perimeter\_m}        & Периметр объекта, м \\
2  & \tbl{compactness}         & Мера компактности: $4\pi S / P^2$ \\
3  & \tbl{elongation}          & Отношение главных полуосей аппроксимирующего эллипса \\
4  & \tbl{convexity}           & Отношение площади объекта к площади его выпуклой оболочки \\
5  & \tbl{mean\_extent\_score} & Средняя вероятность по каналу Extent (выход BoundaryUNet) \\
6  & \tbl{mean\_boundary\_score} & Средняя вероятность по каналу Boundary (выход BoundaryUNet) \\
7  & \tbl{mean\_ndvi}          & Средний NDVI по пикселям объекта \\
8  & \tbl{std\_ndvi}           & Стандартное отклонение NDVI по пикселям объекта \\
9  & \tbl{max\_ndvi}           & Максимальный NDVI по пикселям объекта \\
10 & \tbl{mean\_bsi}           & Средний BSI по пикселям объекта \\
11 & \tbl{mean\_ndwi}          & Средний NDWI по пикселям объекта \\
12 & \tbl{scl\_valid\_frac}    & Доля валидных пикселей по SCL \\
13 & \tbl{edge\_density}       & Плотность рёбер по признаку \tbl{edge\_composite} \\
14 & \tbl{neighbor\_count}     & Число смежных сегментов \\
15 & \tbl{dist\_to\_road\_m}   & Расстояние до ближайшего дорожного объекта, м \\
16 & \tbl{dist\_to\_water\_m}  & Расстояние до ближайшего водного объекта, м \\
17 & \tbl{texture\_entropy}    & Энтропия текстуры по признаку \tbl{ndvi\_entropy} \\
18 & \tbl{aspect\_ratio}       & Отношение ширины к высоте ограничивающего прямоугольника \\
\bottomrule
\end{tabularx}
\end{table}

\subsection{Вспомогательная модель SAM}

Файл \tbl{backend/models/sam\_vit\_h\_4b8939.pth} содержит чекпойнт модели
Segment Anything Model (SAM) архитектуры ViT-H, разработанной Meta AI.
Размер файла составляет приблизительно 2,4~ГБ. Модель применяется
в качестве вспомогательного инструмента интерактивной сегментации
и не задействована в основном автоматическом пайплайне инференса.

% ═══════════════════════════════════════════════════════════════════
\section{Метрики качества}
% ═══════════════════════════════════════════════════════════════════

Оценка качества модели BoundaryUNet~v3 проведена на валидационной выборке
(10~тайлов, 1~690~патчей). Результаты приведены в таблице~\ref{tab:metrics}.

\begin{table}[h!]
\caption{Метрики качества модели BoundaryUNet v3}
\label{tab:metrics}
\begin{tabularx}{\textwidth}{L{110mm}X}
\toprule
\textbf{Метрика} & \textbf{Значение} \\
\midrule
IoU (канал Extent)                              & 0,534 \\
F1-мера (канал Boundary)                        & 0,581 \\
Расстояние Хаусдорфа (95-й перцентиль), пикс.   & 14,27 \\
Максимальное отклонение вывода ONNX от PyTorch  & $< 5{,}1 \cdot 10^{-8}$ \\
\bottomrule
\end{tabularx}
\end{table}

% ═══════════════════════════════════════════════════════════════════
\section{Параметры конфигурации системы}
% ═══════════════════════════════════════════════════════════════════

Параметры инференса задаются посредством переменных окружения,
определённых в файле \tbl{.env} проекта.
Ключевые параметры приведены в таблице~\ref{tab:env}.

\begin{table}[h!]
\caption{Параметры конфигурации системы}
\label{tab:env}
\begin{tabularx}{\textwidth}{L{62mm}L{40mm}X}
\toprule
\textbf{Переменная} & \textbf{Значение по умолчанию} & \textbf{Описание} \\
\midrule
\env{MODE}                     & \tbl{production}    & Режим работы системы \\
\env{ALLOW\_SYNTHETIC\_DATA}   & \tbl{false}         & Разрешение использования синтетических данных \\
\env{ML\_USE\_ONNX}            & \tbl{true}          & Использование ONNX-экспорта при инференсе \\
\env{ML\_MODEL\_PATH}          & \tbl{models/boundary\_unet\_v3\_cpu.onnx} & Путь к файлу модели \\
\env{ML\_MODEL\_NORM\_STATS\_PATH} & \tbl{models/...norm.json} & Путь к файлу статистик нормализации \\
\env{ML\_FEATURE\_PROFILE}     & \tbl{v2\_16ch}      & Идентификатор профиля признаков \\
\env{MODEL\_VERSION}           & \tbl{boundary\_unet\_v3\_cpu} & Версия модели \\
\env{FEATURE\_ML\_PRIMARY}     & \tbl{true}          & ML-признаки как основные \\
\env{USE\_OBJECT\_CLASSIFIER}  & \tbl{true}          & Активация постфильтрации \\
\env{OBJECT\_CLASSIFIER\_PATH} & \tbl{models/object\_classifier.pkl} & Путь к классификатору \\
\env{OBJECT\_MIN\_SCORE}       & \tbl{0.40}          & Порог уверенности классификатора \\
\env{TEMPORAL\_YEARS\_BACK}    & \tbl{1}             & Временной горизонт наблюдений, лет \\
\bottomrule
\end{tabularx}
\end{table}

% ═══════════════════════════════════════════════════════════════════
\section{Методология регионального дообучения}
% ═══════════════════════════════════════════════════════════════════

\subsection{Назначение}

Пайплайн регионального дообучения реализован в скрипте
\tbl{backend/training/train\_regional\_supplement.py} и предназначен
для адаптации базовой модели BoundaryUNet к спектральным и фенологическим
характеристикам конкретных агроклиматических зон Российской Федерации.
Дообучение выполняется на региональных тайлах, загружаемых через
Sentinel Hub API, с последующим переобучением классификатора объектов
и экспортом обновлённых моделей.

\subsection{Параметры регионального профиля}
\label{sec:profiles}

Для каждого региона определяется набор параметров обработки
(\tbl{RegionalProfile}), влияющих как на формирование обучающих данных,
так и на ход дообучения. Параметры профилей для двух реализованных
регионов приведены в таблицах~\ref{tab:krd_profile}
и~\ref{tab:perm_profile}.

\begin{table}[h!]
\caption{Параметры профиля Краснодарского края}
\label{tab:krd_profile}
\begin{tabularx}{\textwidth}{L{50mm}L{24mm}X}
\toprule
\textbf{Параметр} & \textbf{Значение} & \textbf{Обоснование} \\
\midrule
\tbl{min\_cloud\_pct}         & 20\,\%      & Строгий порог: высокая инсоляция юга обеспечивает достаточное число чистых сцен \\
\tbl{temporal\_window\_months}& 12          & Полный годовой цикл для захвата всех фенологических фаз (рис, пшеница, подсолнечник) \\
\tbl{date\_range\_years}      & 2           & Охват двух вегетационных сезонов для учёта межгодовой изменчивости \\
\tbl{min\_field\_coverage}    & 0,05        & Минимальная доля поля в патче (плотная сельскохозяйственная зона) \\
\tbl{learning\_rate}          & $10^{-4}$   & Стандартная скорость дообучения при достаточном объёме данных \\
\tbl{finetune\_epochs}        & 50          & Число эпох дообучения \\
\tbl{loss\_weight\_boundary}  & 2,0         & Повышенный вес: поля имеют чёткие прямолинейные контуры \\
\tbl{loss\_weight\_distance}  & 0,5         & — \\
\tbl{freeze\_encoder\_epochs} & 10          & Первые 10 эпох обучается только декодер — предотвращает деградацию базовых признаков \\
\bottomrule
\end{tabularx}
\end{table}

\begin{table}[h!]
\caption{Параметры профиля Пермского края}
\label{tab:perm_profile}
\begin{tabularx}{\textwidth}{L{50mm}L{24mm}X}
\toprule
\textbf{Параметр} & \textbf{Значение} & \textbf{Обоснование} \\
\midrule
\tbl{min\_cloud\_pct}         & 40\,\%      & Мягкий порог: высокая повторяемость облаков снижает число чистых сцен \\
\tbl{temporal\_window\_months}& 6           & Только вегетационный период (май–октябрь): зима не несёт полезного сигнала \\
\tbl{date\_range\_years}      & 3           & Расширенный диапазон компенсирует дефицит чистых сцен \\
\tbl{min\_field\_coverage}    & 0,03        & Сниженный порог для мелкоконтурных полей (1--10~га) \\
\tbl{learning\_rate}          & $5\cdot10^{-5}$ & Пониженная скорость: небольшой объём данных повышает риск переобучения при высоком LR \\
\tbl{finetune\_epochs}        & 50          & Число эпох дообучения \\
\tbl{augmentation\_strength}  & strong      & Усиленная аугментация компенсирует ограниченность обучающей выборки \\
\tbl{loss\_weight\_boundary}  & 1,5         & Умеренный вес: поля имеют менее чёткие контуры, чем в Краснодарском крае \\
\tbl{freeze\_encoder\_epochs} & 15          & Увеличенный период фиксации энкодера защищает базовые признаки при малом датасете \\
\bottomrule
\end{tabularx}
\end{table}

\subsection{Этапы выполнения пайплайна}

Пайплайн регионального дообучения включает шесть последовательных этапов,
описанных в таблице~\ref{tab:pipeline_stages}.

\begin{table}[h!]
\caption{Этапы пайплайна регионального дообучения}
\label{tab:pipeline_stages}
\begin{tabularx}{\textwidth}{C{6mm}L{46mm}L{54mm}X}
\toprule
\textbf{№} & \textbf{Этап} & \textbf{Описание} & \textbf{Продолжительность} \\
\midrule
1 & Загрузка тайлов            & Получение NPZ-тайлов Sentinel-2 через Sentinel Hub API     & 20--40 мин \\
2 & Формирование меток         & Слабые метки через Overpass API, растеризация, морфология  & 10--20 мин \\
3 & Дообучение UNet            & Региональное дообучение BoundaryUNet с заморозкой энкодера & 1--3 ч \\
4 & Переобучение классификатора& Переобучение HistGBM на региональных объектах-примерах     & ${\sim}15$ мин \\
5 & Экспорт ONNX               & Конвертация обновлённой модели в формат ONNX opset~18      & ${\sim}2$ мин \\
6 & Выгрузка в хранилище       & Загрузка обновлённых файлов в Google Drive через rclone    & ${\sim}5$ мин \\
\bottomrule
\end{tabularx}
\end{table}

\subsection{Запуск пайплайна}

Полный запуск пайплайна для обоих регионов выполняется командой:

\begin{lstlisting}
PYTHONPATH=backend python backend/training/train_regional_supplement.py \
    --regions krasnodar perm \
    --epochs 50 \
    --push-gdrive
\end{lstlisting}

При необходимости пропуска уже выполненных этапов используются флаги
\tbl{{-}{-}skip-fetch} и \tbl{{-}{-}skip-labels}:

\begin{lstlisting}
PYTHONPATH=backend python backend/training/train_regional_supplement.py \
    --regions krasnodar perm \
    --epochs 50 \
    --skip-fetch \
    --skip-labels
\end{lstlisting}

По завершении пайплайна обновляются файлы
\tbl{backend/models/boundary\_unet\_v3\_cpu.pth},
\tbl{backend/models/boundary\_unet\_v3\_cpu.onnx}
и формируется отчёт
\tbl{backend/models/regional\_training\_report\_YYYYMMDD\_HHMMSS.json}.

% ═══════════════════════════════════════════════════════════════════
\section*{Заключение}
\addcontentsline{toc}{section}{Заключение}
% ═══════════════════════════════════════════════════════════════════

Настоящий документ содержит техническое описание набора данных, программных
моделей и методологии обучения, применяемых в программном комплексе AGROSIGNAL.
Описаны методы формирования 16-канального признакового пространства на основе
многовременных снимков Sentinel-2, структура NPZ-тайлов, механизм формирования
слабых меток из данных OpenStreetMap, архитектура модели BoundaryUNet~v3
с многозадачной функцией потерь, а также процедура регионального дообучения
под условия конкретных агроклиматических зон.

Документ подлежит актуализации при выпуске новых версий моделей,
расширении регионального датасета или изменении признакового профиля.

\end{document}
